% GENERATED FILE. Edit canonical lessons, then run npm run generate:books.
% canonical-source-hash: fnv1a64:4d3e6dd13a2b1b7b
% canonical-lessons: IT-C370-maggioranza, IT-C370-minoranza, IT-C370-democrazia, IT-C370-riforma, IT-C370-deputato, IT-R370-ripasso-363-370

\chapter{La maggioranza}
\label{ch:it-la-maggioranza}
\hlchaptermodality{370}

\begin{chapteropening}
\textbf{By the end of this chapter:} \emph{I can talk about the majority and minorities, democracy, a reform and a member of parliament.}

Complete the last lesson of chapter 370: I can talk about the majority and minorities, democracy, a reform and a member of parliament.
\end{chapteropening}

\section[la maggioranza]{la maggioranza — the majority}

\label{lesson:IT-C370-maggioranza}



\begin{quote}
Before the new one: say the Italian for to elect, then the Italian for the voter.
\end{quote}

\subsection*{You'll want to know: la maggioranza}
\textbf{la maggioranza} — \textquotedblleft{}the majority\textquotedblright{}.

\begin{quote}
La maggioranza degli studenti è d'accordo.
\end{quote}

\begin{quote}
\emph{The majority of the students agree.}
\end{quote}

\emph{La maggioranza} of parliament is the governing side.

\begin{cousinweb}
\emph{maggioranza} is built on \emph{maggiore} \textquotedblleft{}greater\textquotedblright{}, inherited from Latin \emph{māior} \textquotedblleft{}greater\textquotedblright{} (Lewis \& Short), the comparative of \emph{magnus} \textquotedblleft{}great\textquotedblright{}.
\end{cousinweb}

\begin{grammarlens}[title={What you've built}]
One more word for telling, explaining and describing at B1.
\end{grammarlens}

\subsection*{Your turn}
\begin{itemize}
\raggedright
  \item \emph{Say it:} \emph{la maggioranza}
  \item \emph{Say it:} \emph{la maggioranza}, once more
  \item \emph{Say it:} a sentence of your own with \emph{la maggioranza}
  \item \emph{Recall:} say the Italian for to elect, then the Italian for the voter, then say \emph{la maggioranza} again
\end{itemize}

\begin{tcolorbox}[breakable,colback=teal!4,colframe=teal!35!black,title={Before you move on}]
What does la maggioranza mean? (\textquotedblleft{}The majority\textquotedblright{}.) Say it once more.
\end{tcolorbox}
\section[la minoranza]{la minoranza — the minority}

\label{lesson:IT-C370-minoranza}



\begin{quote}
Before the new one: say the Italian for the voter, then the Italian for the majority.
\end{quote}

\subsection*{You'll want to know: la minoranza}
\textbf{la minoranza} — \textquotedblleft{}the minority\textquotedblright{}.

\begin{quote}
In questa regione vive una minoranza che parla tedesco.
\end{quote}

\begin{quote}
\emph{A minority who speak German lives in this region.}
\end{quote}

\emph{Essere in minoranza} — to be outnumbered.

\begin{cousinweb}
\emph{minoranza} is built on \emph{minore} \textquotedblleft{}smaller\textquotedblright{}, inherited from Latin \emph{minor} \textquotedblleft{}smaller, less\textquotedblright{} (Lewis \& Short).
\end{cousinweb}

\begin{grammarlens}[title={What you've built}]
One more word for telling, explaining and describing at B1.
\end{grammarlens}

\subsection*{Your turn}
\begin{itemize}
\raggedright
  \item \emph{Say it:} \emph{la minoranza}
  \item \emph{Say it:} \emph{la minoranza}, once more
  \item \emph{Say it:} a sentence of your own with \emph{la minoranza}
  \item \emph{Recall:} say the Italian for the voter, then the Italian for the majority, then say \emph{la minoranza} again
\end{itemize}

\begin{tcolorbox}[breakable,colback=teal!4,colframe=teal!35!black,title={Before you move on}]
What does la minoranza mean? (\textquotedblleft{}The minority\textquotedblright{}.) Say it once more.
\end{tcolorbox}
\section[la democrazia]{la democrazia — the democracy}

\label{lesson:IT-C370-democrazia}



\begin{quote}
Before the new one: say the Italian for the majority, then the Italian for the minority.
\end{quote}

\subsection*{You'll want to know: la democrazia}
\textbf{la democrazia} — \textquotedblleft{}the democracy\textquotedblright{}.

\begin{quote}
In una democrazia tutti possono votare.
\end{quote}

\begin{quote}
\emph{In a democracy everyone can vote.}
\end{quote}

Stress on \emph{-zìa}.

\begin{cousinweb}
\emph{democrazia} is a learned borrowing from Medieval Latin \emph{dēmocratia}, from Greek \emph{dēmokratía} \textquotedblleft{}rule by the people\textquotedblright{} (EtymDB), built from \emph{dêmos} \textquotedblleft{}the people\textquotedblright{} and \emph{krátos} \textquotedblleft{}strength, power\textquotedblright{}.
\end{cousinweb}

\begin{grammarlens}[title={What you've built}]
One more word for telling, explaining and describing at B1.
\end{grammarlens}

\subsection*{Your turn}
\begin{itemize}
\raggedright
  \item \emph{Say it:} \emph{la democrazia}
  \item \emph{Say it:} \emph{la democrazia}, once more
  \item \emph{Say it:} a sentence of your own with \emph{la democrazia}
  \item \emph{Recall:} say the Italian for the majority, then the Italian for the minority, then say \emph{la democrazia} again
\end{itemize}

\begin{tcolorbox}[breakable,colback=teal!4,colframe=teal!35!black,title={Before you move on}]
What does la democrazia mean? (\textquotedblleft{}The democracy\textquotedblright{}.) Say it once more.
\end{tcolorbox}
\section[la riforma]{la riforma — the reform}

\label{lesson:IT-C370-riforma}



\begin{quote}
Before the new one: say the Italian for the minority, then the Italian for the democracy.
\end{quote}

\subsection*{You'll want to know: la riforma}
\textbf{la riforma} — \textquotedblleft{}the reform\textquotedblright{}.

\begin{quote}
La riforma della scuola non piace agli insegnanti.
\end{quote}

\begin{quote}
\emph{Teachers do not like the school reform.}
\end{quote}

The verb is \emph{riformare}.

\begin{cousinweb}
\emph{riforma} is built on \emph{riformare} \textquotedblleft{}to reform\textquotedblright{}, a learned borrowing from Latin \emph{refōrmāre} \textquotedblleft{}to shape again\textquotedblright{} (Lewis \& Short), built on \emph{fōrma} \textquotedblleft{}shape\textquotedblright{}.
\end{cousinweb}

\begin{grammarlens}[title={What you've built}]
One more word for telling, explaining and describing at B1.
\end{grammarlens}

\subsection*{Your turn}
\begin{itemize}
\raggedright
  \item \emph{Say it:} \emph{la riforma}
  \item \emph{Say it:} \emph{la riforma}, once more
  \item \emph{Say it:} a sentence of your own with \emph{la riforma}
  \item \emph{Recall:} say the Italian for the minority, then the Italian for the democracy, then say \emph{la riforma} again
\end{itemize}

\begin{tcolorbox}[breakable,colback=teal!4,colframe=teal!35!black,title={Before you move on}]
What does la riforma mean? (\textquotedblleft{}The reform\textquotedblright{}.) Say it once more.
\end{tcolorbox}
\section[il deputato]{il deputato — the member of parliament, the deputy}

\label{lesson:IT-C370-deputato}



\begin{quote}
Before the new one: say the Italian for the democracy, then the Italian for the reform.
\end{quote}

\subsection*{You'll want to know: il deputato}
\textbf{il deputato} — \textquotedblleft{}the member of parliament, the deputy\textquotedblright{}.

\begin{quote}
Il deputato ha votato contro la riforma.
\end{quote}

\begin{quote}
\emph{The member of parliament voted against the reform.}
\end{quote}

A woman is \emph{la deputata}.

\begin{cousinweb}
\emph{deputato} is a learned borrowing from Late Latin \emph{dēputātus} \textquotedblleft{}assigned, sent\textquotedblright{}, from \emph{dēputāre} \textquotedblleft{}to allot, to assign\textquotedblright{} (Lewis \& Short), built on \emph{putāre} \textquotedblleft{}to reckon\textquotedblright{}.
\end{cousinweb}

\begin{grammarlens}[title={What you've built}]
One more word for telling, explaining and describing at B1.
\end{grammarlens}

\subsection*{Your turn}
\begin{itemize}
\raggedright
  \item \emph{Say it:} \emph{il deputato}
  \item \emph{Say it:} \emph{il deputato}, once more
  \item \emph{Say it:} a sentence of your own with \emph{il deputato}
  \item \emph{Recall:} say the Italian for the democracy, then the Italian for the reform, then say \emph{il deputato} again
\end{itemize}

\begin{tcolorbox}[breakable,colback=teal!4,colframe=teal!35!black,title={Before you move on}]
What does il deputato mean? (\textquotedblleft{}The member of parliament, the deputy\textquotedblright{}.) Say it once more.
\end{tcolorbox}
\section[(dialogue)]{Ripasso dei capitoli 363–370}

\label{lesson:IT-R370-ripasso-363-370}



\begin{quote}
Say the Italian for the citizen and the citizenship.
\end{quote}

\subsection*{Your turn}
Say these aloud:

\begin{itemize}
\raggedright
  \item the citizen, the citizenship, the right, the freedom, the equality
  \item the protest, to protest, the demonstration, to demonstrate, the banner
  \item the justice, unfair, the injustice, the prejudice, the stereotype
  \item the poverty, the unemployment, the discrimination, to discriminate, the immigration
  \item the voluntary work, the volunteer, the association, the charity, the community
  \item the living together, the tolerance, the respect, to respect, the integration
  \item the political party, the parliament, the minister, to elect, the voter
  \item the majority, the minority, the democracy, the reform, the member of parliament
\end{itemize}

\begin{tcolorbox}[breakable,colback=teal!4,colframe=teal!35!black,title={Before you move on}]
The citizen? (\textbf{il cittadino}.) The citizenship? (\textbf{la cittadinanza}.) The right? (\textbf{il diritto}.) The freedom? (\textbf{la libertà}.) The equality? (\textbf{l'uguaglianza}.) The protest? (\textbf{la protesta}.) To protest? (\textbf{protestare}.) The demonstration? (\textbf{la manifestazione}.) To demonstrate? (\textbf{manifestare}.) The banner? (\textbf{lo striscione}.) The justice? (\textbf{la giustizia}.) Unfair? (\textbf{ingiusto}.) The injustice? (\textbf{l'ingiustizia}.) The prejudice? (\textbf{il pregiudizio}.) The stereotype? (\textbf{lo stereotipo}.) The poverty? (\textbf{la povertà}.) The unemployment? (\textbf{la disoccupazione}.) The discrimination? (\textbf{la discriminazione}.) To discriminate? (\textbf{discriminare}.) The immigration? (\textbf{l'immigrazione}.) The voluntary work? (\textbf{il volontariato}.) The volunteer? (\textbf{il volontario}.) The association? (\textbf{l'associazione}.) The charity? (\textbf{la beneficenza}.) The community? (\textbf{la comunità}.) The living together? (\textbf{la convivenza}.) The tolerance? (\textbf{la tolleranza}.) The respect? (\textbf{il rispetto}.) To respect? (\textbf{rispettare}.) The integration? (\textbf{l'integrazione}.) The political party? (\textbf{il partito}.) The parliament? (\textbf{il parlamento}.) The minister? (\textbf{il ministro}.) To elect? (\textbf{eleggere}.) The voter? (\textbf{l'elettore}.) The majority? (\textbf{la maggioranza}.) The minority? (\textbf{la minoranza}.) The democracy? (\textbf{la democrazia}.) The reform? (\textbf{la riforma}.) The member of parliament? (\textbf{il deputato}.)
\end{tcolorbox}
